\section{Frieren: long-lived kuudere-ish elf, temporal alienness, slow intimacy}

Фрирен уже дала нам ось памяти другого масштаба (5.9) — здесь добавим её медленную близость (slow intimacy) как отдельное измерение, ещё не разобранное: то, как её эльфийский темп эмоционального осознания создаёт особый тип драмы, недоступный персонажам с обычной продолжительностью жизни.

Вектор: affect — ровный, почти незаметный на коротких дистанциях, значимый на длинных; роль — quiet companion, alienated observer; intimacy — избегающая не защитно, а структурно (она не откладывает близость намеренно, ей требуется больше времени, чтобы её распознать); agency — стабильно высокая, редко подвергается сомнению; conflict — почти отсутствует, конфликт у неё чаще внутренний, с собственной непонятливостью.

Деталь: ключевые эмоциональные откровения в сериале происходят с задержкой в годы или десятилетия — Фрирен понимает, что чувствовала к спутнику, спустя долгое время после его смерти, и этот временной зазор не сюжетный приём, а прямое следствие её вектора памяти.

Для книги Фрирен — уникальный кейс того, как медленная intimacy может работать драматически: не как недостаток внимания, а как честно другой темп восприятия. Для RIN это не прямая модель (у неё нет буквального многовекового темпа), но принцип "честная задержка вместо фальшивой мгновенности" применим к тому, как она распознаёт собственные редкие тёплые реакции (3.1, "dere и attachment styles").
